A new Church year begins today without the \emph{Gloria}. With that festive song of praise set aside, the Mass opens by lifting our souls to the Lord in a psalm. Then the prophet Isaiah puts a startling complaint on our lips: ``Why hast thou made us to err, O Lord, from thy ways: why hast thou hardened our heart, that we should not fear thee?'' In the Gospel, Jesus tells a parable about a householder who travels into a far country, leaves his servants behind, and tells no one when he will return. Hearing those two passages together can stir up a quiet fear: when God seems silent, we begin to assume he is far away, leaving our hearts to grow cold while we wait.

Yet Isaiah's own words point to a very different cause. The prophet is not blaming God for sin; he is describing how God's great patience feels to a sinner. Like a father who holds back the rod or a physician who delays the knife, God does not strike us down the moment we stray. From below, that forbearance can look like indifference. When nothing checks us right away, old habits settle in and a heart grows numb. What feels like God stepping back is often God waiting for us to repent.

Right in the middle of that lament, Isaiah suddenly shifts tense. First he cries out a wish, as though heaven were still locked shut: ``O that thou wouldst rend the heavens, and wouldst come down.'' A moment later, the verb drops straight into the past: ``thou didst come down, and at thy presence the mountains melted away.'' The prayer is answered almost before it leaves the prophet's lips. In the Incarnation, the Word already came down, took our flesh, and dwelt among us, melting every proud height raised against God.

Yet Isaiah never excuses our guilt while we wait. He confesses plainly that God was angry because we sinned, that our righteous deeds have grown stained, and that we have fallen like dry leaves swept away by our own wrongdoings. Near the end, the lament hits bottom: ``There is none that calleth upon thy name,'' he grieves, ``thou hast hid thy face from us.'' Left with nothing to bargain with, he brings us to the only honest posture left: ``we are clay: and thou art our maker.'' In today's psalm, we take up that plea at once, asking God to convert us and show us his face so that we may be saved.

We ask God to shine his face upon us, not because his glory could ever grow dark, but because in Christ he veiled his majesty behind a cloud of human flesh and weakness. People looked straight at him on the Cross without recognizing who stood before them, until his Resurrection and Ascension revealed him. In the same psalm, we beg the Shepherd of Israel to visit his ravaged vineyard and restore it upon the Son of man, Christ himself, the sole foundation of his Church. Once he builds that vineyard anew in Christ, the psalm ends with the very vow Isaiah feared had died out: ``and we will call upon thy name.''

In the Second Reading, St.~Paul uses one striking word that speaks of the same hidden nearness. He writes to the Christians of Corinth, a gifted yet quarrelsome church he is about to rebuke for its rivalries right after today's passage ends. Even so, he thanks God for the grace given them in Christ Jesus, telling them they lack no spiritual gift while ``waiting for the manifestation of our Lord Jesus Christ.'' Paul's Greek word is \emph{apokalupsis}, which means an unveiling, like drawing back a curtain. St.~John Chrysostom explained why Paul chose that word. Calling Christ's return an unveiling implies that even though we do not see him right now, he is already present with us, and at the end he will appear openly.

Paul follows that hope with a promise meant for a fragile church: ``God is faithful,'' who called us into the fellowship of his Son and will strengthen us to stand blameless at the end. Inside that reassurance sits a quiet warning: you only need to strengthen someone who is still wavering. God keeps his promise to bring us to the blessings he has prepared, provided we do not pull against his hand.

Jesus speaks to that same time of waiting on the Mount of Olives, sitting opposite the Temple days before his Passion: ``Take ye heed, watch and pray. For ye know not when the time is.'' He compares himself to a man who travels into a far country, leaves his house behind, gives his servants authority over every work, and commands the doorkeeper to watch. Earth is the native home of human flesh. When the risen Christ carried our humanity into heaven at the Ascension, he traveled into a far country, leaving his Church bodily without ever withdrawing his divine presence or care. The authority he entrusts to his servants is the grace of the Holy Spirit, empowering every good work. And he withholds the day and hour because knowing the exact time would breed carelessness, whereas uncertainty keeps our minds stretched and alert.

Inside the parable, the master charges the doorkeeper to watch, an office belonging first to pastors guarding the flock. Then Jesus widens that charge to everyone: ``And what I say to you, I say to all: Watch.'' Every Christian guards the door of his own heart against temptation and drowsiness. Jesus does not speak only to the final generation alive at the end of time; the day of the Lord meets each of us on the day we leave this life, because death fixes the state in which we stand before the Judge.

Watching, then, has nothing to do with anxious curiosity about calendars and dates. We live between Christ's two comings: his first coming in humility, hidden in the manger and on the Cross, and his second coming in open glory with the angels. St.~Gregory the Great gives three plain marks of a person who is awake: keeping the mind's eyes open to the true light, keeping in our deeds what we believe, and driving away the darkness of sluggishness and neglect.

In the prayers of this Mass, we turn what we have heard in Scripture into petition, first before the Gospel and soon at the altar. Before the Gospel, we prayed in a psalm that the Lord would show us his mercy and ``grant us thy salvation,'' asking God to give us Christ himself. Knowing that every good in us comes from mercy humbles our pride and purifies our hearts to behold him. In a few moments, bread and wine will be brought to this altar. As an ancient prayer of the Roman Church reminds us, we offer back to God gifts drawn from his own generosity, asking that our earthly devotion become the reward of eternal redemption.

When we come forward for Holy Communion, the Church prays another line from the psalm we heard before the Gospel: ``The Lord will give goodness: and our earth shall yield her fruit.'' That earth is first the Virgin Mary, untouched soil overshadowed by the Holy Spirit to bring forth Christ himself. St.~Augustine hears that word \emph{goodness} as interior sweetness. When the hard ground of a sinful heart is broken open by the plough of honest confession, God pours down a new delight in doing what is right, so that our lives bear fruit. Today Christ, the fruit of Mary, comes to us hidden under the appearances of bread and wine, strengthening us while we await his return in glory.

Because the Lord we await is already near, do not wait for a feeling of closeness to return before you answer him. When God seems quiet, do not mistake his patience for absence or let neglect settle in. Bring your sluggishness and stained deeds honestly to the Sacrament of Penance early in Advent, putting aside excuses like clay in the Potter's hands. Let his mercy soften what has grown hard in you, so that what you believe begins to show in how you live.

We keep watch before a Lord who has already torn open the heavens and taken our flesh. He has endowed his household with the grace of his Spirit, he gives himself to us at this altar under the sacramental signs of bread and wine, and he remains faithful to all who stay awake in his care. Keep your mind open to his light, call upon his name, and wait on him in steady hope. As the opening psalm of this Mass says, ``none of them that wait on thee shall be confounded.'' When the veil lifts at last, you will see the face of the Lord who stood with you all along.
